%!TEX root = ../template.tex
%%%%%%%%%%%%%%%%%%%%%%%%%%%%%%%%%%%%%%%%%%%%%%%%%%%%%%%%%%%%%%%%%%%%
%% chapter4.tex
%% NOVA thesis document file
%%
%% Chapter with lots of dummy text
%%%%%%%%%%%%%%%%%%%%%%%%%%%%%%%%%%%%%%%%%%%%%%%%%%%%%%%%%%%%%%%%%%%%

\typeout{NT FILE chapter4.tex}%

\chapter{Related Work}
\label{cha:related-work}

% Quick intro to the chapter
In this chapter, we will review some of the related work in the field of image classification
of aerial/satellite images using machine learning techniques. This review will help us understand
the current landscape of research in this area and identify potential work
that can be used as a reference for our own research.

\section{Remote Sensing Image Classification}

The article \textit{Remote Sensing Image Scene Classification Meets Deep Learning} \cite{cheng2020remote} provides a systematic review of deep learning methods for remote
sensing image scene classification, covering more than 160 papers. It compares the accuracies of
several models using the UC-MERCED dataset \cite{yang2010bag}, AID dataset \cite{xia2017aid}, and
NWPU-RESISC45 dataset \cite{cheng2017remote}, offering a comprehensive overview of their performance.

The UC-Merced dataset, contains 21 scene classes with 100
land-use images per class, resulting in a total of 2100 images. Each image has a fixed resolution
of 256x256 pixels with a spatial resolution of 0.3 meters.

The AID dataset, consists of 30 scene classes,
each containing between 220 and 420 images. These images, totaling 10,000, were cropped from Google
Earth imagery and fixed at 600x600 pixels. The dataset is multisourced and multiresolution, with
pixel resolutions ranging from about 8 meters to 0.5 meters.

The NWPU-RESISC45 dataset, is
the largest among the three, with 45 scene categories and 700 images per category, amounting
to 31,500 images in total. These images were sourced from Google Earth, fixed at 256x256 pixels,
and span over 100 countries and regions. The pixel resolution varies from about 30 meters to
0.2 meters for most scene categories. Notably, the NWPU-RESISC45 is the most challenging of the
three datasets and therefore carries more weight in our evaluation of model performance.

Although our specific problem focuses on fire risk classification rather than scene
classification, this article still provides a valuable overview of state-of-the-art
methods in the field of remote sensing image classification.
\subsection{Potential Challenges}

We can specifically relate to the challenge of interclass similarity. It may be difficult to
distinguish between close risk classes, such as low and moderate risk, due to overlapping vegetation
characteristics. These classes can be visually similar, making them hard to differentiate even for
the human eye.

Additionally imaging conditions can also have a significant impact, as illustrated in the Figures \ref{fig:vc83}
and \ref{fig:vc102}, taken from the \gls{dgt} open API. These orthophotos were taken at different times of the day, resulting in varying shadows
and lighting conditions. Such variations can affect the appearance of vegetation and surrounding
areas, further complicating the accurate classification of fire risk.

\begin{figure}[ht]
    \centering
    \begin{minipage}[b]{0.45\textwidth}
        \centering
        \includegraphics[width=\textwidth]{../Figures/VC83.png}
        \caption{Very High Risk House}
        \label{fig:vc83}
    \end{minipage}
    \hfill
    \begin{minipage}[b]{0.45\textwidth}
        \centering
        \includegraphics[width=\textwidth]{../Figures/VR95.png}
        \caption{High Risk House}
        \label{fig:vc102}
    \end{minipage}
    \caption{Very High and High Fire Risk Areas respectively in \gls{irg} format.}
    \label{fig:veryhigh-vs-high}
\end{figure}

To address the challenges posed by shadows, we can explore using satellite images from open API
sources to capture house photos at more optimal times of the day. However, this approach may
not be feasible due to the resolution and quality limitations of freely available images
compared to the \gls{dgt} orthophotos, unless we invest in higher-quality data. Therefore, we will
consider this as a last resort if shadows prove to be a significant issue. Ideally, it
would be highly beneficial to develop a model that is robust to these shadow variations,
as this could also provide extra information about the height of the vegetation.

\subsection{Deep learning-based scene classification algorithms}

The article \cite{cheng2020remote} divides deep learning-based scene classification
algorithms into three categories: the autoenconder-based methods, CNN-based methods,
and GAN-based methods.

\subsubsection{Autoencoder-based methods}
Autoencoders have been widely used in earlier years for remote sensing image scene classification due
to their ability to automatically learn mid-level visual representations from unlabeled data.
While autoencoder-based methods have achieved good results in remote sensing image scene
classification, most of them do not fully exploit scene class information, which limits their
ability to learn the best discriminative features for distinguishing different scene classes.
Therefore, since their performance is generally lower than that of CNN-based methods as proved in
\cite{cheng2020remote}, we will not focus on autoencoder-based methods in this work.

\subsubsection{CNN-based methods}
The use of \gls{cnn} for remote sensing image classification tasks has become increasingly
popular in recent years. The CNN-based methods of remote sensing image classification
can be divided into three categories: using pretrained \gls{cnn} as feature extractors, fine-tuning
pretrained \gls{cnn} on target datasets, and training \gls{cnn} from scratch.

\smallskip

Using \textbf{pretrained \gls{cnn} as feature extractors} and then fusing or combining the features
extracted by existing \gls{cnn} is a common practice as referred in \cite{cheng2020remote}. It
is simple and effective on small-scale datasets. Some of the best performing and
noteworthy referred are the MSCP (Multilayer Stacked Covariance Pooling) \cite{MSCP} and
the FACNN (Feature Aggregation Convolution Neural Network) \cite{FACNN}, with accuracies
of 98.36\% and 98.81\% respectively for the UC-MERCED dataset.

The MSCP uses a pretrained model to extract multilayer feature maps, those feature maps are
stacked together and a covariance matrix is calculated. Each entry of the resulting covariance
matrix provides a natural and innovative way to exploit complementary information provided
by different layers feature maps. Finally, the extracted covariance matrices are used as features
for classification by a Support Vector Machine (SVM).

The FACNN, on the other hand, comprises three main components: a backbone pipeline, a
convolutional feature encoding module, and a softmax classifier as seen in Figure \ref{fig:facnn-arch}.

\begin{figure}[ht]
    \centering
    \includegraphics[width=1\textwidth]{../Figures/FACNN-arch.png}
    \caption{FACNN architecture \cite{FACNN}.}
    \label{fig:facnn-arch}
\end{figure}
The convolutional feature encoding module is designed to encode intermediate convolutional features into a unified
convolutional representation through end-to-end training. Initially, the intermediate
convolutional features are fed to a pooling layer to standardize the width and height of
different features to 7 x 7. Subsequently, a concatenation layer is used to combine these
features by their channels. This is followed by 1 x 1 convolutional layers with ReLU activation
functions to fuse the channel-wise information of the different features. A global average
pooling layer is then used to produce the convolutional representation. This representation
is merged with the second fully connected (FC) feature to create a discriminative scene
representation, which is ultimately fed to a softmax classifier to derive the semantic
labels of remote sensing scenes.

Both MSCP and FACNN have achieved good results, therefore, if we decide to use pretrained
\gls{cnn} as feature extractors, these models could be a good starting point.

\smallskip

\textbf{Fine-tuning pretrained \gls{cnn}} is particularly useful when the amount of training samples
is insufficient to train a new CNN from scratch. The idea is to fine-tune an already trained
CNN on target datasets while optimizing its loss functions. Among the best-performing works
compared in \cite{cheng2020remote} using the UC-MERCED dataset are the SF-CNN (Scale-Free
Convolutional Neural Network) \cite{SF-CNN} and the GBN (Gated Bidirectional Network) \cite{GBN},
with accuracies of 98.57\% and 99.05\%, respectively. However, for the NWPU-RESISC45 dataset, models
such as Hydra \cite{hydra}, which employs an ensemble of \gls{cnn}, proved to be superior. Hydra achieved
an accuracy of 94.51\%, the highest among the models tested on this dataset, compared to the SF-CNN's
accuracy of 92.55\%. Therefore, an ensemble approach like Hydra might be a solution to consider for our
problem.

The hydra framework makes use of two CNN popular architectures ResNet (residual networks)
\cite{he2016deep} and DenseNet (dense convolutional networks) \cite{huang2017densely}.
The training process consisted on creating the body of the Hydra by optimizing the chosen
architectures using weights pretrained on ImageNet \cite{deng2009imagenet} as a starting
point. Then with the obtained weights, they are copied to several \gls{cnn}, which are further
optimized to form the heads of the Hydra. To achieve a final classification, an image
is fed to all Hydra's heads and their outputs are fused to generate a single decision as
seen in Figure \ref{fig:hydra-testing}.

\begin{figure}[ht]
    \centering
    \includegraphics[width=0.8\textwidth]{../Figures/hydra-testing.png}
    \caption{Hydra framework testing flowchart from \cite{hydra}.}
    \label{fig:hydra-testing}
\end{figure}

Even though fine-tuning pretrained \gls{cnn} achieved remarkable performances, this type of approach has some limitations,
such as learned features not always being entirely suitable for the characteristics of target
datasets, which can be an inconvenience, needing to modify pretrained \gls{cnn}.

\smallskip

\textbf{Training \gls{cnn} from scratch} has obtained astonishing results as referred in \cite{cheng2020remote}.
However they require numerous annotated samples to fine-tune already trained \gls{cnn} or train a network
from scratch which is not only time-consuming but also computationally expensive. One of the
best-performing models using this approach is the SCCov (Skip-Connected Covariance) \cite{sccov}.
The SSCov model adds skip connections and covariance pooling to \gls{cnn}, reducing the amount
of parameters and achieving a good classification performance with a performance of 99.05\% for
the UC-MERCED dataset, the same as the SF-CNN model.

For our problem at hand, training \gls{cnn} from scratch might not be the best approach due to the
size of our dataset.

\subsubsection{GAN-based methods}
The key method for unsupervised learning, Generative Adversarial Networks (GANs), contrasts with
CNN-based methods in the context of remote sensing, which typically require massive labeled samples
to train models effectively, presenting a significant challenge due to the labor-intensive nature
of sample annotation.

In response to this challenge, some studies have started employing GANs for scene classification
\cite{attgan}, \cite{martagan}. GANs offer a distinct advantage by enabling the generation of
high-quality synthetic samples that can be used to augment training datasets, thereby alleviating
the dependency on large volumes of labeled data.

Despite their potential, GAN-based scene classification methods face several limitations. Compared to
CNN-based methods, GAN-based approaches are relatively underexplored in the literature, and their
performance is generally considered inferior. Moreover, most GAN-based methods cannot be trained
end-to-end as they often require labels for training an additional classifier. This requirement
diminishes the efficiency of GANs as a stand-alone solution for scene classification.

\subsection{Conclusions}
There are many different architectures and techniques that have provided good results
for the image scene classification problem. The best approaches by far were the ones
that followed CNN-based methods. Both the autoencoder-based methods and the GAN-based
methods have shown worse results while also being less explored in the literature.

As we do not have a huge and varied dataset, training a CNN from scratch is definitely
not the best approach for our problem. Instead, we will try both fine-tuning an already
pretrained CNN and using \gls{cnn} as feature extractors to find out which one fits best
our problem. Both the Hydra and the FACNN architectures are good starting points.

Moreover, some architectures, such as the FACNN, may require our images to be in a specific
format, necessitating a reduction in image resolution. While this might lead to a loss of
information, it can also result in faster training times. On the other hand, the Hydra
architecture, which uses an ensemble of ResNet and DenseNet models, offers flexibility
in handling different image sizes, allowing us to maintain the original resolution of the
images, even if at a higher computational cost.
